% M10 load-balance study: what the ocean-phase imbalance is, and what balancing it costs.
% Build:  module load texlive/live2025-gcc-13.3.0 && make
\documentclass[10pt,a4paper]{article}

\usepackage[utf8]{inputenc}
\usepackage[T1]{fontenc}
\usepackage{lmodern}
\usepackage[margin=17mm,top=14mm,bottom=12mm]{geometry}
\usepackage{amsmath,amssymb}
\usepackage{booktabs}
\usepackage{microtype}
\usepackage{parskip}
\setlength{\parskip}{2pt plus 1pt}
\setlength{\tabcolsep}{5pt}
\pagestyle{empty}

\begin{document}

{\large\bfseries Load imbalance in FESOM2/Kokkos: its origin, and the cost of removing it}\\[1pt]
{\small 6 August 2026 --- Levante, 128-core CPU nodes. Per-rank per-phase wall clocks;
MPI wait measured inside MPI by PMPI interposition.}

\textbf{The problem.} At production rank counts the model spends about half of each step
waiting inside MPI (51\,\% on CORE2 at 864 ranks, 53\,\% on fArc at 2048). The largest block
of real computation is the ocean phase, and its per-rank cost varies by a factor of 4.8. A
bulk-synchronous phase costs its slowest rank rather than its mean, so this spread is a direct
tax on the step: 5.2\,ms of a 48.2\,ms step on CORE2 (11\,\%) and 15.9\,ms of 84.7\,ms on fArc
(19\,\%), comparable to the entire barotropic solve.

\textbf{What the imbalance is.} Correlating per-rank ocean compute time against per-rank mesh
content over all 2048 fArc ranks identifies it unambiguously:

\begin{center}\vspace{-4pt}\footnotesize
\begin{tabular}{lr}
\toprule
correlation with per-rank ocean compute time & $r$ \\
\midrule
3D nodes owned ($\sum$ levels over owned columns) & \textbf{0.967} \\
2D nodes owned & 0.003 \\
solver wait vs.\ 3D nodes owned & $-0.771$ \\
\bottomrule
\end{tabular}\vspace{-4pt}
\end{center}

\noindent
The imbalance is vertical. The partition balances \emph{surface} nodes almost exactly (310--313
per rank, a spread of 1.01), but the work is per water column, and 3D nodes span
\textbf{1550 to 14571, a factor of 9.4}. Ocean compute follows as
$2.24\,\mathrm{ms}$ per 1000 3D nodes, which reproduces the measured 8.8--42.5\,ms range. The
negative correlation for solver wait completes the picture: deep-column ranks wait least
because they are the stragglers, and the shallow ranks idle at the solve's first reduction.
The imbalance is bathymetry, not sea ice.

Vertical balance depends on how the partition was weighted, and it varies across the meshes in
use: NG5 and DARS are dual-weighted (3D spread 1.04--1.05), fArc and CORE2 are not (9.40, 9.60).

\textbf{The test.} CORE2 partitions were regenerated at 864 ranks under 2D-only and dual
2D+3D vertex weighting from a single partitioner binary, and run in one allocation with one
model binary, so the decomposition is the only variable. All arms are byte-identical on the
mesh definition; only the partition files differ.

\begin{center}\vspace{-4pt}\footnotesize
\begin{tabular}{lrrrr}
\toprule
partition & 3D balance & halo (nodes/rank) & s/step & vs.\ shipped \\
\midrule
shipped & 9.60 & 42 & 0.0440 & --- \\
2D-only weighting & 9.32 & 42 & 0.0443 & $+0.68$\,\% \\
dual 2D+3D weighting & \textbf{1.05} & \textbf{59} & 0.0482 & $\mathbf{+9.55}$\,\% \\
\bottomrule
\end{tabular}\vspace{-4pt}
\end{center}

\noindent
The 2D-only arm reproduces the shipped partition's imbalance and runtime, so the partitioner
contributes nothing and the 9.55\,\% is the weighting alone.

\textbf{The result.} Dual weighting removes the imbalance and still costs more than it saves:

\begin{center}\vspace{-4pt}\footnotesize
\begin{tabular}{lrrr}
\toprule
ms/step, 864 ranks & 2D-only & dual & $\Delta$ \\
\midrule
ocean compute, min / mean / max & 3.8 / 13.0 / 19.4 & 10.1 / 14.7 / 17.5 & spread $5.1\to1.73$ \\
ocean imbalance tax (max $-$ mean) & 6.4 & 2.8 & $-3.6$ \\
ocean wait & 6.7 & 4.9 & $-1.8$ \\
total compute & 22.9 & 25.6 & $+2.7$ \\
solver wait & 9.0 & 10.2 & $+1.2$ \\
ice-dynamics wait & 4.0 & 5.1 & $+1.1$ \\
\midrule
total & 45.2 & 48.1 & $\mathbf{+2.9}$ \\
\bottomrule
\end{tabular}\vspace{-4pt}
\end{center}

\noindent
The balancing behaves as intended --- previously idle ranks take on real work and ocean wait
drops 1.8\,ms --- but balancing the vertical requires cutting more edges, and the halo grows
40\,\% (42 to 59 nodes per rank). That adds 2.7\,ms of compute and raises wait in the two
communication-heavy phases, together exceeding the 1.8\,ms recovered.

\textbf{Conclusion.} The imbalance is real, bathymetric, and worth 6.4\,ms of a 45\,ms step,
but it is not recoverable by vertex weighting: the halo penalty exceeds the imbalance gain.
Only the two endpoints of the edgecut-against-balance trade are measured, and the dual
weighting already softens its 3D criterion, so an intermediate weight remains unexplored.

\end{document}
